\documentclass[10pt,a4paper]{amsart}
\usepackage[margin=19mm]{geometry}
\usepackage{amsmath,amssymb,mathtools}
\usepackage[scr=rsfs,cal=euler]{mathalfa}
\usepackage{xfrac}
\usepackage[unicode,hidelinks,bookmarks=false]{hyperref}
\newcommand{\ie}{i.e.\ }
\newcommand{\cf}{cf.\ }
\newcommand{\resp}{respectively\ }
\newcommand{\Subsection}{\subsection}
\providecommand{\nobreakdash}{\nobreak\hbox{-}}
\setlength{\emergencystretch}{2em}
\multlinegap0pt
\usepackage{float}
\usepackage{xcolor}
\usepackage[shortlabels]{enumitem}
\newtheorem{lemma}{Lemma}
\newtheorem{proposition}{Proposition}
\title{Stress-Testing Neural Network Verifiers with Provably Robust Instances}
\author{David Troxell, Yulia Alexandr, Sofia Hunt, Stephanie Lei, Guido Mont\'ufar}
\date{Source: arXiv:2605.17153v1 (CC BY 4.0)}
\begin{document}
\maketitle

\noindent\emph{Source and licence.} Stress-Testing Neural Network Verifiers with Provably Robust Instances. Version arXiv:2605.17153v1 (CC BY 4.0), distributed by arXiv under the Creative Commons Attribution 4.0 International licence, http://creativecommons.org/licenses/by/4.0/, as recorded in the arXiv licence metadata for this exact version and shown on the abstract page https://arxiv.org/abs/2605.17153v1. Full licence notice: \texttt{licenses/arxiv-2605.17153-CC-BY-4.0.txt}.\par\smallskip

\noindent\textit{Editorial scope.} Counted unit: the article's proposition prop:linear\_attention\_perturbation (source0.tex lines 1098--1112). The article's complete original proof of it is delivered with it (source0.tex lines 1114--1181, one \texttt{proof} environment of 68 lines). No mathematics is written, repaired or re-derived: the statement and the proof are the article's own lines, in the article's own notation, and no retained line is substituted or re-flowed.  Retained uncounted as the source-local context the proof needs: lines 260--265; lines 1015--1032; lines 1089--1096.  Nothing was omitted from any retained span, so the package carries no omission record.  No retained line contains a citation, so the wrapper carries no bibliography and no bibliography entry.  No retained line contains a figure environment, an image-inclusion command or drawing code, so no image is delivered and the package declares no asset dependency.  Editorial formatting only, in addition: an AMS \texttt{amsart} preamble that restates the article's own declarations and macros used by the retained lines, namely the environments \texttt{lemma}, \texttt{proposition} on separate counters, together with the standard packages those lines need.  The retained environments print in this document as Lemma 1, Proposition 1 in order of appearance, whereas the article prints the same items with the numbering of its own full document.  The complete native source of the article as distributed in the arXiv e-print for this exact version is bundled unchanged as \texttt{sources/200773/source0.tex}.
\begin{equation}
    w_{i j_i^\star}(X)
    \geq
    \rho_i \sum_{j\neq j_i^\star} w_{ij}(X).
    \label{eq:linear_attention_dominance}
\end{equation}

\begin{lemma}[Closeness to the dominant value]
\label{lem:linear_dominant_closeness}
If row $i$ satisfies the dominance condition in \eqref{eq:linear_attention_dominance} throughout $\mathcal B_\epsilon(X_0)$, then for every $X\in\mathcal B_\epsilon(X_0)$,
\begin{equation}
    \left\|
    \mathrm{Attn}_{\mathrm{lin}}(X)[i,:]
    -
    V(X)[j_i^\star,:]
    \right\|_2
    \leq
    \frac{1}{1+\rho_i}
    \max_{j\neq j_i^\star}
    \left\|
    V(X)[j,:]-V(X)[j_i^\star,:]
    \right\|_2.
    \label{eq:app_linear_closeness}
\end{equation}
\end{lemma}

\begin{equation}
    \|V(X)[j,:]-V_0[j,:]\|_2
    =
    \|\delta^{(j)}W_V\|_2
    \leq
    \epsilon L_V.
    \label{eq:app_value_drift}
\end{equation}

\begin{proposition}[Dominant-key attention-output perturbation]
\label{prop:linear_attention_perturbation}
Suppose every row satisfies the dominance condition in \eqref{eq:linear_attention_dominance} throughout $\mathcal B_\epsilon(X_0)$, and let $\rho=\min\limits_i\rho_i$. Then for every $X\in\mathcal B_\epsilon(X_0)$,
\begin{equation}
    \|\mathrm{Attn}_{\mathrm{lin}}(X)-\mathrm{Attn}_{\mathrm{lin}}(X_0)\|_{2,\infty}
    \leq
    \frac{2}{1+\rho}D_V(X_0)
    +
    \left(
        1+\frac{2}{1+\rho}
    \right)
    \epsilon L_V.
    \label{eq:app_linear_attention_bound}
\end{equation}
\end{proposition}

\begin{proof}
Fix a row $i$. Insert the value vector associated with the dominant key at both $X$ and $X_0$:
\begin{equation}
\begin{aligned}
    &\left\|
    \mathrm{Attn}_{\mathrm{lin}}(X)[i,:]
    -
    \mathrm{Attn}_{\mathrm{lin}}(X_0)[i,:]
    \right\|_2 \\
    &\quad\leq
    \left\|
    \mathrm{Attn}_{\mathrm{lin}}(X)[i,:]-V(X)[j_i^\star,:]
    \right\|_2
    +
    \left\|
    V(X)[j_i^\star,:]-V_0[j_i^\star,:]
    \right\|_2 \\
    &\qquad+
    \left\|
    V_0[j_i^\star,:]-\mathrm{Attn}_{\mathrm{lin}}(X_0)[i,:]
    \right\|_2.
\end{aligned}
\label{eq:app_linear_three_terms}
\end{equation}
By Lemma~\ref{lem:linear_dominant_closeness} applied at $X$,
\begin{equation}
    \left\|
    \mathrm{Attn}_{\mathrm{lin}}(X)[i,:]-V(X)[j_i^\star,:]
    \right\|_2
    \leq
    \frac{1}{1+\rho}
    \max_{j\neq j_i^\star}
    \|V(X)[j,:]-V(X)[j_i^\star,:]\|_2.
    \label{eq:app_linear_term_A}
\end{equation}
For every $j\neq j_i^\star$, adding and subtracting nominal values gives
\begin{equation}
\begin{aligned}
    \|V(X)[j,:]-V(X)[j_i^\star,:]\|_2
    &\leq
    \|V_0[j,:]-V_0[j_i^\star,:]\|_2 \\
    &\quad+
    \|V(X)[j,:]-V_0[j,:]\|_2
    +
    \|V(X)[j_i^\star,:]-V_0[j_i^\star,:]\|_2 \\
    &\leq
    D_V(X_0)
    +
    2\epsilon L_V.
\end{aligned}
\label{eq:app_linear_value_spread_perturbed}
\end{equation}
Thus the first term in \eqref{eq:app_linear_three_terms} is bounded by $\frac{D_V(X_0)+2\epsilon L_V}{(1+\rho)}$. The middle term is bounded by $\epsilon L_V$ using \eqref{eq:app_value_drift}. The last term is bounded by $\frac{D_V(X_0)}{(1+\rho)}$ by Lemma~\ref{lem:linear_dominant_closeness} applied at $X_0$. Summing the three bounds gives
\begin{equation}
    \left\|
    \mathrm{Attn}_{\mathrm{lin}}(X)[i,:]-
    \mathrm{Attn}_{\mathrm{lin}}(X_0)[i,:]
    \right\|_2
    \leq
    \frac{2}{1+\rho}D_V(X_0)
    +
    \left(
        1+\frac{2}{1+\rho}
    \right)
    \epsilon L_V.
\end{equation}
Taking the maximum over rows gives the $\|\cdot\|_{2,\infty}$ bound.
\end{proof}

\end{document}
